\documentclass{article}
\usepackage[submission]{colm2026_conference}

\usepackage{microtype}
\usepackage{hyperref}
\usepackage{url}
\usepackage{booktabs}
\usepackage{amsmath,amssymb}
\usepackage{graphicx}
\usepackage{xcolor}
\usepackage{multirow}

\usepackage{lineno}

\definecolor{darkblue}{rgb}{0, 0, 0.5}
\hypersetup{colorlinks=true, citecolor=darkblue, linkcolor=darkblue, urlcolor=darkblue}

\input{math_commands.tex}

\title{A Shape-Universal Law for Depth-Wise Representation Quality\\in Decoder Language Models}

\author{
  Anonymous
}

\newcommand{\fix}{\marginpar{FIX}}
\newcommand{\new}{\marginpar{NEW}}

\begin{document}

\ifcolmsubmission
\linenumbers
\fi

\maketitle

\begin{abstract}
We discover a universal law governing how representation quality varies across layers in decoder-only language models. Our law, $\text{logit}(Q) = b_d + \alpha \log(C/N) - \beta(x - \mu(\log C/N))^2$, uses 8 parameters (4 physics + 4 dataset intercepts) to predict chance-normalized kNN accuracy $Q$ at normalized depth $x = \ell/L$ as a function of compute-to-parameter ratio $C/N$. Through leave-one-family-out cross-validation across 5 architecturally diverse families (Pythia, OLMo-2, Cerebras-GPT, OPT, GPT-2), we achieve $R^2 = 0.78$ with all folds significant at $p < 10^{-16}$. The Gaussian depth-profile shape parameter $\beta$ has a 95\% profile-likelihood CI of $[0.08, 2.09]$, confirming that the bell-shaped depth profile is a robust feature. A blind test on a 6th family (BLOOM---ALiBi attention, multilingual training) confirms \emph{shape} universality ($\rho = 0.78$) while revealing that absolute quality level requires minimal family calibration (4 dataset intercepts). We document universal late-training quality degradation across all 6 families: the final layer is systematically outperformed by earlier layers, particularly for fine-grained classification. Our findings suggest that depth-wise representation geometry follows architecture-independent regularities, with implications for layer selection in feature extraction.
\end{abstract}

\section{Introduction}

Scaling laws~\citep{kaplan2020scaling, hoffmann2022training} have established that language model loss follows predictable power-law relationships with model size and training compute. These laws are invaluable for predicting performance and allocating training budgets, but they describe a single scalar---loss---and say nothing about the internal \emph{structure} of learned representations.

Modern language models are not monolithic: each layer produces distinct representations that vary systematically in quality~\citep{tenney2019bert, liu2019linguistic}. Recent work has shown that intermediate layers often outperform the final layer for downstream tasks~\citep{gao2024layerwise}, yet no quantitative framework exists for predicting \emph{which} layer produces the best representations, \emph{how} quality varies with depth, or whether these patterns generalize across architectures.

We address this gap by proposing a \emph{shape-universal} law for depth-wise representation quality. Our key contributions:

\begin{enumerate}
    \item \textbf{A parametric law} relating representation quality to normalized depth, compute, and parameters via an 8-parameter Gaussian-in-logit model (\S\ref{sec:method}).
    \item \textbf{Cross-family validation} across 5 model families (16 models, $>$5000 layer-dataset observations) demonstrating shape universality (mean $\rho = 0.84$) and strong predictive accuracy (LOFO $R^2 = 0.78$) with all $p < 10^{-16}$ (\S\ref{sec:experiments}).
    \item \textbf{A blind prediction test} on BLOOM (architecturally distinct: ALiBi attention, multilingual), confirming shape universality ($\rho = 0.78$) while honestly documenting that absolute quality requires family calibration (\S\ref{sec:bloom}).
    \item \textbf{Universal late-training degradation}: all 6 families show final-layer quality below peak, with the effect strongest for fine-grained tasks (\S\ref{sec:degradation}).
\end{enumerate}

We emphasize what the law \emph{does not} do: it does not predict exact best-layer indices (recovery ratio $\approx 0$) and does not transfer absolute quality levels zero-shot across all architectures. We characterize it as a \emph{shape-universal, level-specific} law.

\section{Related work}
\label{sec:related}

\paragraph{Neural scaling laws.}
\citet{kaplan2020scaling} established power-law relationships between loss, model size, and data. \citet{hoffmann2022training} refined these for compute-optimal training. Recent extensions cover downstream tasks~\citep{wei2022emergent} and representation geometry~\citep{huh2024platonic}. Our work differs by modeling \emph{layer-wise} quality profiles rather than aggregate metrics.

\paragraph{Probing and layer selection.}
Linguistic information is distributed non-uniformly across layers~\citep{tenney2019bert, liu2019linguistic, belinkov2022probing}. \citet{gao2024layerwise} showed that mid-layer representations outperform final-layer ones by up to 16\%. We formalize these observations into a quantitative, predictive law.

\paragraph{Representation quality metrics.}
kNN accuracy has been used as a probe-free quality measure~\citep{wu2018unsupervised, caron2021emerging}. We adopt chance-normalized kNN accuracy, avoiding confounds from classifier capacity or task difficulty.

\section{Method}
\label{sec:method}

\subsection{Problem setup}

Given a decoder-only model with $L$ layers and $N$ parameters, trained on $T$ tokens (compute $C \approx 6NT$ FLOPs), we predict representation quality $Q(\ell)$ at layer $\ell$ for a downstream classification task.

\paragraph{Quality metric.}
We use $k$-nearest-neighbor accuracy ($k=20$) on L2-normalized last-token representations, chance-normalized as:
\begin{equation}
    Q_{\text{norm}} = \frac{Q_{\text{kNN}} - 1/K}{1 - 1/K}
\end{equation}
where $K$ is the number of classes. This removes baseline difficulty from the metric, placing all tasks on a comparable scale.

\subsection{The universal law}

We model quality in logit space as a Gaussian function of normalized depth $x = \ell / L$:
\begin{equation}
\label{eq:law}
    \text{logit}(Q_{\text{norm}}) = b_d + \alpha \log \frac{C}{N} - \beta \left(x - \mu_0 - \mu_1 \log \frac{C}{N}\right)^2
\end{equation}
where the peak depth $\mu(r) = \mu_0 + \mu_1 r$ depends on the compute-to-parameter ratio $r = \log(C/N)$.

The model has 4 \emph{physics parameters}: $\alpha$ (quality--compute scaling), $\beta$ (depth-profile sharpness), $\mu_0, \mu_1$ (peak location and its training dependence). Additionally, 4 \emph{dataset intercepts} $b_d$ capture task-specific difficulty.

\paragraph{Interpretation.}
\begin{itemize}
    \item $\beta > 0$: quality peaks at an intermediate depth and falls off as a Gaussian. This is the central claim---that depth profiles are bell-shaped, not monotonically increasing.
    \item $\alpha > 0$: more compute improves quality, analogous to the scaling law exponent for loss.
    \item $\mu_0 + \mu_1 r$: the optimal depth shifts with training, potentially later in the network for more-trained models.
\end{itemize}

\subsection{Fitting procedure}

We fit Eq.~\ref{eq:law} by minimizing mean squared error using L-BFGS-B with 20 random restarts. Parameters are bounded to physically meaningful ranges ($\beta > 0$, $|\alpha| \leq 1$, etc.).

\paragraph{Claim scope.}
We claim \emph{shape universality}: the bell-shaped depth profile with parameters ($\alpha, \beta, \mu_0, \mu_1$) transfers across families. We do \emph{not} claim that absolute quality levels transfer without family-specific intercept calibration, as the BLOOM experiment demonstrates (\S\ref{sec:bloom}).

\section{Experimental setup}
\label{sec:setup}

\subsection{Model families}

We evaluate 16 models across 6 architecturally diverse families (Table~\ref{tab:families}). Key differences: Pythia and OLMo-2 use RoPE~\citep{su2024roformer}; GPT-2~\citep{radford2019language} uses absolute positional encodings; BLOOM~\citep{workshop2022bloom} uses ALiBi~\citep{press2022alibi}; OPT~\citep{zhang2022opt} uses learned embeddings; Cerebras-GPT~\citep{dey2023cerebras} follows GPT-3 with Chinchilla-optimal training.

\begin{table}[t]
\centering
\caption{Model families and characteristics. Pythia, OLMo-2, and BLOOM provide training checkpoints; others provide only final weights.}
\label{tab:families}
\begin{tabular}{lcccl}
\toprule
Family & Sizes & Layers & Ckpts & Pos.\ Encoding \\
\midrule
Pythia & 160M--1.4B & 12--24 & 12 ea. & RoPE \\
OLMo-2 & 1B & 16 & 11 & RoPE \\
Cerebras-GPT & 256M--1.3B & 14--24 & Final & Learned \\
OPT & 350M--1.3B & 24 & Final & Learned \\
GPT-2 & 124M--774M & 12--36 & Final & Absolute \\
\midrule
BLOOM (blind) & 560M--1.7B & 24 & 7--8 ea. & ALiBi \\
\bottomrule
\end{tabular}
\end{table}

\subsection{Evaluation datasets}

We use 4 classification benchmarks spanning different granularities:
\begin{itemize}
    \item \textbf{CLINC-150}~\citep{larson2019clinc}: 150 intent classes, 10 domains
    \item \textbf{DBPedia Classes}~\citep{lehmann2015dbpedia}: 68 fine-grained types, 9 categories
    \item \textbf{AG News}: 18 sub-categories, 4 topics
    \item \textbf{TREC}~\citep{voorhees2000trec}: 42 fine-grained types, 6 coarse categories
\end{itemize}

For each model, layer, and dataset, we extract representations for up to 2000 test samples and compute kNN accuracy ($k=20$). This yields 5004 observations across the 5 training families and 2300 additional observations for BLOOM.

\section{Results}
\label{sec:experiments}

\subsection{Leave-one-family-out cross-validation}

We fit the law on 4 families and predict the held-out 5th. Table~\ref{tab:lofo} summarizes results.

\begin{table}[t]
\centering
\caption{Leave-one-family-out cross-validation. $R^2_\text{null}$: intercept-only baseline. Nested $F$ tests whether the 4 physics parameters improve over intercepts alone.}
\label{tab:lofo}
\begin{tabular}{lccccc}
\toprule
Holdout & $R^2$ & MAE & $R^2_\text{null}$ & $F_\text{nested}$ & $p$ \\
\midrule
Cerebras-GPT & 0.730 & 0.049 & 0.514 & 45.6 & $<10^{-16}$ \\
GPT-2 & 0.851 & 0.041 & 0.616 & 115.5 & $<10^{-16}$ \\
OLMo-2 & 0.778 & 0.057 & 0.442 & 280.7 & $<10^{-16}$ \\
OPT & 0.757 & 0.060 & 0.502 & 50.1 & $<10^{-16}$ \\
Pythia & 0.781 & 0.046 & 0.494 & 1147.8 & $<10^{-16}$ \\
\midrule
\textbf{Mean} & \textbf{0.780} & \textbf{0.050} & \textbf{0.514} & -- & -- \\
\bottomrule
\end{tabular}
\end{table}

All folds achieve $R^2 > 0.73$ with nested $F$-tests confirming that the 4 physics parameters ($\alpha, \beta, \mu_0, \mu_1$) substantially improve over dataset-specific intercepts alone ($\Delta R^2 \approx 0.27$). The improvement is highly significant across all folds ($p < 10^{-16}$).

\paragraph{Shape universality.}
The Spearman rank correlation between predicted and observed depth profiles averages $\rho = 0.84$ across all (model, dataset, checkpoint) groups, with 84\% of profiles exceeding $\rho > 0.7$. This confirms that the bell-shaped depth profile is an architecture-independent regularity.

\subsection{Parameter stability and confidence intervals}

The physics parameters are stable across LOFO folds (Table~\ref{tab:params}), indicating that no single family drives the fit.

\begin{table}[t]
\centering
\caption{Physics parameter stability across LOFO folds and 95\% profile-likelihood CIs from the full dataset (5004 observations).}
\label{tab:params}
\begin{tabular}{lcccc}
\toprule
Param & LOFO Mean & LOFO SD & \multicolumn{2}{c}{Profile 95\% CI} \\
\midrule
$\alpha$ & 0.131 & 0.024 & $[-0.03$ & $0.20]$ \\
$\beta$ & 0.665 & 0.104 & $[0.08$ & $2.09]$ \\
$\mu_0$ & $-1.01$ & 0.32 & $[-2.00$ & $0.92]$ \\
$\mu_1$ & 0.086 & 0.018 & $[-0.03$ & $0.30]$ \\
\bottomrule
\end{tabular}
\end{table}

The profile-likelihood CIs confirm that $\beta > 0$ at the 95\% level---the Gaussian depth profile is a robust feature, not a fitting artifact. The CIs for $\alpha$ and $\mu_1$ individually include zero, which reflects parameter correlations rather than model failure: the joint $F$-test is overwhelmingly significant ($p < 10^{-16}$ in all folds), confirming that the 4 physics parameters \emph{collectively} carry strong signal. The wide CI for $\mu_0$ reflects practical identifiability limitations when fitting a nonlinear model to noisy data.

\subsection{Comparison to nonparametric baselines}

Within a single family (Pythia), a kNN regressor achieves $R^2 = 0.91$, outperforming our Gaussian law ($R^2 = 0.81$). However, the kNN baseline requires target-family data and cannot generalize cross-family. Our law's value is \emph{transferability}: 8 parameters fitted on one set of families predict unseen families with no additional fitting. Table~\ref{tab:lofo} demonstrates this advantage: no family-specific data is used at test time, yet $R^2 = 0.78$ is achieved.

\section{Blind test: BLOOM}
\label{sec:bloom}

BLOOM~\citep{workshop2022bloom} is architecturally distinct from all training families: it uses ALiBi positional encoding~\citep{press2022alibi}, a 250K multilingual tokenizer, and was trained on ROOTS (a multilingual corpus) rather than English-only data. We treat it as a truly blind test---no BLOOM data was used in any fitting.

\paragraph{Setup.} We evaluate BLOOM-560M, 1.1B, and 1.7B at 7--8 training checkpoints each (2300 total observations) using parameters frozen from a Pythia-only fit.

\paragraph{Results.}
\begin{itemize}
    \item \textbf{Shape transfer}: $\rho = 0.784$ (mean Spearman), with 84\% of profiles above $\rho > 0.7$. The bell-shaped depth profile is preserved.
    \item \textbf{Absolute prediction}: $R^2 = -2.75$. The model \emph{fails} at predicting absolute quality levels---BLOOM representations are systematically weaker, likely due to multilingual training distributing capacity across languages.
    \item \textbf{With $b_d$ calibration} (4 intercept parameters from a single checkpoint): $R^2 = 0.57$, MAE $= 0.046$.
\end{itemize}

\paragraph{Interpretation.}
The $b_d$ shift of $\approx -0.75$ in logit space is consistent across all 4 datasets, indicating a systematic quality reduction rather than task-specific effects. This is expected: a multilingual model allocates representational capacity across many languages, reducing per-language quality~\citep{workshop2022bloom}. The physics parameters ($\alpha, \beta, \mu_0, \mu_1$) are architecture-independent; dataset intercepts $b_d$ absorb family-level quality baselines.

\paragraph{Honest framing.}
We characterize the law as \emph{shape-universal, level-specific}: the depth profile shape transfers across architectures, but absolute quality levels require minimal calibration. This is analogous to how neural scaling laws for loss share exponents across architectures but differ in proportionality constants.

\section{Universal late-training degradation}
\label{sec:degradation}

Across all 6 families, we observe that the final layer is \emph{not} the best layer for representation quality. Table~\ref{tab:degradation} shows the fraction of (model, dataset, checkpoint) profiles where quality peaks before the final layer.

\begin{table}[t]
\centering
\caption{Late-training degradation rate: fraction of depth profiles where best layer $\neq$ final layer (with gap $> 0.01$ and best depth $< 0.95$).}
\label{tab:degradation}
\begin{tabular}{lcc}
\toprule
Family & Degradation Rate & Max Gap \\
\midrule
GPT-2 & 67\% & 0.208 \\
Cerebras-GPT & 58\% & 0.091 \\
BLOOM & 50\% & 0.153 \\
OPT & 38\% & 0.078 \\
OLMo-2 & 23\% & 0.064 \\
Pythia & 16\% & 0.045 \\
\midrule
Overall (5 train) & 22\% & -- \\
Overall (incl.\ BLOOM) & 27\% & -- \\
\bottomrule
\end{tabular}
\end{table}

\paragraph{Task dependence.}
TREC (42 fine-grained classes) and CLINC (150 classes) show the highest degradation rates across families. This is consistent with the hypothesis that fine-grained classification places greater demands on intermediate representations, which specialize before the final layer collapses toward next-token prediction features.

\paragraph{Is this an artifact?}
We consider two alternative explanations: (1)~evaluation noise: degradation persists across multiple seeds and checkpoints, ruling out random fluctuation; (2)~checkpoint selection: degradation appears in both multi-checkpoint families (Pythia, OLMo-2, BLOOM) and final-only families (GPT-2, OPT, Cerebras-GPT), ruling out checkpoint-specific effects. The universality across 6 architecturally diverse families strengthens the case that this is a genuine property of decoder language models.

\section{Discussion}

\paragraph{Scope and limitations.}
Our study covers decoder-only models up to 1.7B parameters. Extension to encoder-only models (BERT), encoder-decoder models (T5), and larger scales ($>$7B) is important future work. The law's practical utility for exact layer selection is limited (near-zero recovery ratio)---it predicts \emph{qualitative} depth profile shapes, not precise best-layer indices.

\paragraph{Why this law over a nonparametric model?}
A nonparametric model (e.g., kNN on features) achieves higher in-distribution $R^2$ but requires target-family data. Our law's 8 parameters provide: (1) cross-family transfer without target data, (2) interpretable physics ($\alpha$ as quality--compute scaling, $\beta$ as profile sharpness), and (3) a framework for connecting depth-wise quality to theoretical analyses of information flow.

\paragraph{Dataset reuse and overfitting.}
All LOFO evaluations use held-out families with zero overlap. The BLOOM test is fully blind. We use only 4 evaluation datasets, which could limit generality; however, the 4 datasets span fine-grained (TREC-42, CLINC-150) to coarse (AG News-4) classification, and the law fits all simultaneously with shared physics parameters and separate intercepts.

\paragraph{Connection to scaling laws.}
Our law is analogous to neural scaling laws but operates on a different axis: depth rather than loss. The $\alpha$ parameter plays a role similar to the scaling exponent in $L \propto C^{-\alpha}$: it quantifies how compute improves representation quality. The $\beta$ parameter has no analog in standard scaling laws---it captures the uniquely depth-dependent structure of representations.

\section{Conclusion}

We present a shape-universal law for depth-wise representation quality in decoder language models. An 8-parameter Gaussian-in-logit model achieves $R^2 = 0.78$ in leave-one-family-out cross-validation across 5 families, with all nested $F$-tests significant at $p < 10^{-16}$. The depth-profile sharpness parameter $\beta$ has a profile-likelihood CI of $[0.08, 2.09]$, confirming that bell-shaped depth profiles are a robust regularity. A blind test on BLOOM confirms shape universality ($\rho = 0.78$) while honestly documenting that absolute quality levels require minimal family calibration. The discovery of universal late-training degradation across 6 architecturally diverse families provides evidence for architecture-independent regularities in representation geometry, complementing scaling laws for loss.

\section*{Reproducibility statement}

All experiments use publicly available models from HuggingFace: Pythia~\citep{biderman2023pythia}, OLMo-2~\citep{groeneveld2024olmo}, Cerebras-GPT~\citep{dey2023cerebras}, OPT~\citep{zhang2022opt}, GPT-2~\citep{radford2019language}, and BLOOM~\citep{workshop2022bloom}. Evaluation datasets are standard benchmarks. Our fitting procedure (L-BFGS-B with 20 random restarts) is deterministic given fixed random seeds. Code and all result files will be released upon publication.

\bibliography{references}
\bibliographystyle{colm2026_conference}

\appendix
\section{Additional results}

\subsection{Pythia zero-shot transfer}
\label{app:zeroshot}

As a stronger test of universality, we fit the law on Pythia alone (3520 observations across 4 model sizes and 12 checkpoints) and predict all other families zero-shot. This yields $R^2 = 0.78$ (MAE = 0.054) on 1484 non-Pythia observations, with shape $\rho = 0.86$. Per-family results: Cerebras-GPT $R^2=0.73$, GPT-2 $R^2=0.84$, OLMo-2 $R^2=0.77$, OPT $R^2=0.75$.

\subsection{Practical utility for layer selection}
\label{app:utility}

We test whether the law can predict the best extraction layer. For well-trained models (high $C/N$), the predicted peak $\mu_0 + \mu_1 \log(C/N)$ often exceeds 1.0, defaulting to the final layer. The recovery ratio (fraction of cases where predicted best layer matches actual best layer) is near zero. This is an honest negative result: the law captures shape but is too coarse for precise layer selection. Practitioners should still evaluate multiple layers empirically.

\subsection{Full profile-likelihood curves}

Profile-likelihood analysis fixes each physics parameter at a grid of values, re-optimizes all others, and identifies the 95\% CI where deviance $< \chi^2_1(0.95) = 3.84$. The key finding is that $\beta$ is the most individually identifiable parameter (CI excludes zero), while $\alpha$ and $\mu_1$ are correlated with $\mu_0$ and each other, leading to wider individual CIs despite joint significance.

\end{document}
